\documentclass{article}
\usepackage[utf8]{inputenc}
\usepackage{a4wide}

\begin{document}

\section*{สำรวจภูเขา}

กำหนดอาร์เรย์\texttt{arr}ของเลขจำนวนเต็มมาให้ ให้สำรวจหาภูเขาจากอาร์เรย์ย่อย

(subarray) ที่มีเงื่อนไขดังต่อไปนี้



1. ความยาวอาร์เรย์ย่อยไม่น้อยกว่า 3



2. มีดัชนี\texttt{i}(ดัชนีเริ่มต้นจาก 0) โดยที่ 0 \textless{} i \textless{}

ความยาวอาร์เรย์ -- 1 และ



\texttt{arr{[}0{]}\ \textless{}\ arr{[}1{]}\ \textless{}\ ...\ \textless{}\ arr{[}i\ –\ 1{]}\ \textless{}\ arr{[}i{]}}และ



\texttt{arr{[}i{]}\ \textgreater{}\ arr{[}i\ +\ 1{]}\ \textgreater{}\ ...\ \textgreater{}\ arr{[}arr.length\ –\ 1{]}}



อาร์เรย์ของจำนวนเต็ม ให้หาความยาวของ\ul{\textbf{อาร์เรย์ย่อยภูเขา}}ที่ยาวที่สุด ตอบ 0

หากไม่มีอาร์เรย์ย่อยภูเขานั้น



\hfill\break



\hfill\break



{\def\LTcaptype{none} % do not increment counter

\begin{longtable}[]{@{}

  >{\raggedright\arraybackslash}p{(\linewidth - 4\tabcolsep) * \real{0.3333}}

  >{\raggedright\arraybackslash}p{(\linewidth - 4\tabcolsep) * \real{0.3333}}

  >{\raggedright\arraybackslash}p{(\linewidth - 4\tabcolsep) * \real{0.3333}}@{}}

\toprule\noalign{}

\begin{minipage}[b]{\linewidth}\centering

ข้อมูลเข้า

\end{minipage} & \begin{minipage}[b]{\linewidth}\centering

ข้อมูลออก

\end{minipage} & \begin{minipage}[b]{\linewidth}\centering

คำอธิบาย

\end{minipage} \\

\midrule\noalign{}

\endhead

\bottomrule\noalign{}

\endlastfoot

\begin{minipage}[t]{\linewidth}\raggedright

6\\

2 1 4 3 2 5\\

\strut

\end{minipage} & 5 & \begin{minipage}[t]{\linewidth}\raggedright

4~ -\/-\textgreater~ 1 4 3 2\\

\strut

\end{minipage} \\

\begin{minipage}[t]{\linewidth}\raggedright

3\\

2 2 2\\

\strut

\end{minipage} & 0 & \begin{minipage}[t]{\linewidth}\raggedright

0~ -\/-\textgreater~ No mountain\\

\strut

\end{minipage} \\

\begin{minipage}[t]{\linewidth}\raggedright

11\\

6 -2 3 5 7 9 2 0 -2 -1 -5\\

\strut

\end{minipage} & 8 & \begin{minipage}[t]{\linewidth}\raggedright

8~ -\/-\textgreater~ -2 3 5 7 9 2 0 -2\\

\strut

\end{minipage} \\

\begin{minipage}[t]{\linewidth}\raggedright

14\\

100 18 45 11 36 56 56 57 72 54 10 57 83 45\\

\strut

\end{minipage} & 5 & \begin{minipage}[t]{\linewidth}\raggedright

5~ -\/-\textgreater~ 56 57 72 54 10 ฐาน (56) ซ้ำ ใช้ได้แค่ตัวเดียว\\

\strut

\end{minipage} \\

\begin{minipage}[t]{\linewidth}\raggedright

6\\

3 5 7 7 5 2\\

\strut

\end{minipage} & 0 & \begin{minipage}[t]{\linewidth}\raggedright

0~ -\/-\textgreater~ 7 ซ้ำกันสองตัวไม่ได้เพราะภูเขาต้องยอดแหลม\\

\strut

\end{minipage} \\

\end{longtable}

}



\subsection*{Input}
บรรทัดแรก จำนวนเต็ม\texttt{N}เป็นขนาดของอาร์เรย์\texttt{arr}



บรรทัดที่สอง เป็นชุดเลขจำนวนเต็มแทนความสูง ณ แต่ละจุด แต่ละค่าจะเว้นด้วยช่องว่าง



\subsection*{Output}
จำนวนเต็มแทนความยาว\ul{\textbf{สูงสุด}}ของอาร์เรย์ย่อยภูเขา\\



\end{document}
